\qquad A seguir apresenta-se a função responsável por orquestrar, para
cada instância, a geração do vetor, a execução do algoritmo de
ordenação selecionado (recebido como um ponteiro de função, o que
permite reutilizar a mesma rotina para os cinco algoritmos deste
trabalho) com a respectiva medição de tempo, e o registro dos
arquivos de entrada, saída e tempo dentro da pasta correspondente ao
algoritmo --- pasta essa que também é criada automaticamente pelo
próprio programa, caso ainda não exista (arquivos
\texttt{algoritmos.h} e \texttt{main.cpp}). Por brevidade, foram omitidas
as mensagens de erro exibidas quando alguma pasta ou arquivo não pode
ser criado. No caso do Merge Sort, o tempo medido inclui a alocação do
vetor auxiliar, por ela fazer parte do algoritmo.

{\small
\begin{verbatim}
bool criarPasta(const string &caminho) {
#if defined(_WIN32)
    int resultado = _mkdir(caminho.c_str());
#else
    int resultado = mkdir(caminho.c_str(), 0777);
#endif
    return resultado == 0 || errno == EEXIST;
}

bool criarPastas(const string &caminho) {
    for (size_t i = 0; i <= caminho.size(); i++) {
        if (i == caminho.size() || caminho[i] == '/') {
            if (i > 0 && !criarPasta(caminho.substr(0, i)))
                return false;
        }
    }
    return true;
}

bool criarEstruturaPastas(const string &algoritmo) {
    const char *grupos[3] = {"Arquivos de Entrada",
                             "Arquivos de Saida",
                             "Arquivos de Tempo"};
    const char *categorias[3] = {"Crescente", "Decrescente", "Random"};
    for (int g = 0; g < 3; g++) {
        for (int c = 0; c < 3; c++) {
            string caminho = algoritmo + "/" + grupos[g] +
                             "/" + categorias[c];
            if (!criarPastas(caminho)) return false;
        }
    }
    return true;
}

bool operacoes(const string &algoritmo, FuncaoOrdenacao funcao,
               int tipo, int tamanho) {
    if (!criarEstruturaPastas(algoritmo)) return false;

    clock_t start_t, end_t;
    double tempoGasto;

    int *vetor = gerarSequencia(tipo, tamanho);

    if (!salvarEntrada(algoritmo, tipo, tamanho, vetor)) {
        delete[] vetor;
        return false;
    }

    start_t = clock();
    funcao(vetor, tamanho);
    end_t = clock();

    tempoGasto = (end_t - start_t) / (double)CLOCKS_PER_SEC;

    bool gravou = salvarTempo(algoritmo, tipo, tamanho, tempoGasto) &&
                  salvarSaida(algoritmo, tipo, tamanho, vetor);

    delete[] vetor;
    return gravou;
}
\end{verbatim}
}
